% ============================================================
%  packages.tex  —  Barcha paketlar shu yerda e'lon qilinadi
% ============================================================

% --- Kodlash va shrift ---
\usepackage[T1]{fontenc}
\usepackage[utf8]{inputenc}
\usepackage{lmodern}
\usepackage{microtype}

% --- Sahifa o'lchamlari ---
\usepackage[a4paper, top=2.5cm, bottom=2.5cm, left=3cm, right=2.5cm]{geometry}

% --- Matematika ---
\usepackage{amsmath}
\usepackage{amssymb}
\usepackage{amsthm}
\usepackage{mathtools}
\usepackage{bm}

% --- Teorema muhitlari ---
\usepackage{thmtools}

% --- Grafiklar ---
\usepackage{tikz}
\usepackage{pgfplots}
\pgfplotsset{compat=1.18}
\usetikzlibrary{
  arrows.meta,
  angles,
  quotes,
  calc,
  patterns,
  patterns.meta,
  decorations.pathmorphing,
  decorations.markings,
  decorations.pathreplacing,
  intersections,
  backgrounds,
  fit,
  shapes.geometric,
  math
}
\usepackage{subcaption}

% --- Jadvallar ---
\usepackage{booktabs}
\usepackage{array}
\usepackage{tabularx}

% --- Ranglar ---
\usepackage{xcolor}

% --- Qutichalar ---
\usepackage[most]{tcolorbox}
\tcbuselibrary{theorems, skins, breakable}

% --- Ro'yxatlar ---
\usepackage{enumitem}

% --- Havolalar ---
\usepackage[
  colorlinks = true,
  linkcolor  = blue!55!black,
  citecolor  = green!45!black,
  urlcolor   = blue!65!black
]{hyperref}

% --- Adabiyotlar ---
\usepackage[backend=biber, style=alphabetic, sorting=nyt, maxbibnames=99]{biblatex}
\addbibresource{bib/references.bib}

% --- Boshqalar ---
\usepackage{setspace}
\usepackage{caption}
\captionsetup{font=small, labelfont=bf}
